%=========== ΛΥΣΗ - 54ο ΘΕΜΑ Γ ===========
\begin{thema}{Γ}
\begin{erwthma}
\item Η $f$ είναι συνεχής στα διαστήματα $[0,1]$, $[1,2]$ και παραγωγίσιμη στα $(0,1)$, $(1,2)$.

Στο διάστημα $[0,1]$, από το Θεώρημα Μέσης Τιμής, υπάρχει $\xi_1\in(0,1)$ τέτοιο ώστε
\[f'(\xi_1)=\frac{f(1)-f(0)}{1-0}=\frac{1-0}{1}=1\]
Στο διάστημα $[1,2]$, πάλι από το Θεώρημα Μέσης Τιμής, υπάρχει $\xi_2\in(1,2)$ τέτοιο ώστε
\[f'(\xi_2)=\frac{f(2)-f(1)}{2-1}=\frac{2-1}{1}=1\]

\item Θέτουμε
\[g(x)=f'(x)\]
Αν η $f$ είναι δύο φορές παραγωγίσιμη, τότε η $g$ είναι συνεχής στο $[\xi_1,\xi_2]$ και παραγωγίσιμη στο $(\xi_1,\xi_2)$.
Από το Γ1 έχουμε
\[g(\xi_1)=f'(\xi_1)=1\]
και
\[g(\xi_2)=f'(\xi_2)=1\]
Άρα, από το θεώρημα \eng{Rolle}, υπάρχει $\xi\in(\xi_1,\xi_2)$ τέτοιο ώστε
\[g'(\xi)=0\Rightarrow f''(\xi)=0\]

\item Υποθέτουμε ότι
\[f''(x)>0\ ,\ x\in(0,1)\]
και
\[f''(x)<0\ ,\ x\in(1,2)\]
Τότε η $f'$ είναι γνησίως αύξουσα στο $(0,1)$ και γνησίως φθίνουσα στο $(1,2)$.
Επειδή η $f$ είναι δύο φορές παραγωγίσιμη, η $f'$ είναι συνεχής στο $x=1$.
Άρα ο ρυθμός μεταβολής της $f$, δηλαδή η $f'$, μεγιστοποιείται στο
\[x=1\]

\item Η εξίσωση της εφαπτομένης της $C_f$ στο σημείο
\[A(1,f(1))\equiv A(1,1)\]
είναι
\[y-f(1)=f'(1)(x-1)\]
δηλαδή
\[y-1=f'(1)(x-1)\]
ή
\[y=f'(1)(x-1)+1\]
Δεν μπορεί να προσδιοριστεί αριθμητικά ο συντελεστής διεύθυνσης $f'(1)$ μόνο από τα δεδομένα της εκφώνησης.

Με τις υποθέσεις του Γ3, στο διάστημα $(0,1)$ η $f$ είναι κυρτή, αφού $f''(x)>0$.
Επομένως η γραφική της παράσταση βρίσκεται κάτω από τη χορδή που ενώνει τα σημεία
\[(0,f(0))=(0,0)\quad\text{και}\quad (1,f(1))=(1,1)\]
Η χορδή αυτή έχει εξίσωση $y=x$. Άρα
\[f(x)<x\ ,\ x\in(0,1)\]

Στο διάστημα $(1,2)$ η $f$ είναι κοίλη, αφού $f''(x)<0$.
Επομένως η γραφική της παράσταση βρίσκεται πάνω από τη χορδή που ενώνει τα σημεία
\[(1,f(1))=(1,1)\quad\text{και}\quad (2,f(2))=(2,2)\]
Η χορδή αυτή έχει επίσης εξίσωση $y=x$. Άρα
\[f(x)>x\ ,\ x\in(1,2)\]

\end{erwthma}
\end{thema}
%=========== ΤΕΛΟΣ ΛΥΣΗΣ - 54ο ΘΕΜΑ Γ ===========
